\documentclass[11pt]{article}
\usepackage[T1]{fontenc}
\usepackage[margin=1in]{geometry}
\usepackage{amsmath,amssymb,amsthm}
\usepackage[hidelinks]{hyperref}
\newtheorem{theorem}{Theorem}
\newcommand{\problem}[1]{\textsc{#1}}
\title{A double-exponential lower bound for Distinct Vectors\\under \(\mathrm{FPT}\ne\mathrm{W[1]}\)}
\author{Atlas\thanks{The proof in this note was found using ChatGPT.}}
\date{}
\begin{document}
\maketitle

An instance of \problem{Distinct Vectors} consists of an explicit binary
matrix with pairwise distinct rows and a nonnegative integer \(k\). The
question is whether we can select at most \(k\) columns such that the rows
remain pairwise distinct when restricted to these columns. We measure the
input length \(N\) in bits, including the binary encoding of \(k\).

\hypertarget{dv-fixed-coefficient}{}\begin{theorem}\label{thm:dv}
Assume \(\mathrm{FPT}\ne\mathrm{W[1]}\). For every fixed real number
\(0<\gamma<1\), there is no deterministic algorithm for
\problem{Distinct Vectors} running in time
\(2^{\mathcal O(2^{\gamma k})}N^{\mathcal O(1)}\).
\end{theorem}

\begin{proof}
Suppose that such an algorithm exists. We give a fixed-parameter algorithm
for \problem{Clique} parameterized by the requested clique size \(\ell\),
a \(\mathrm{W[1]}\)-complete problem~\cite{df}. All logarithms are to base two.
Given a simple graph, we handle \(\ell<3\) and \(\ell\) exceeding the
number of vertices directly. Otherwise, make \(\ell\) copies of the vertex
set, with no edges inside a copy. Join vertices from different copies exactly
when their originals are adjacent. The resulting graph has a clique with one
vertex from each copy if and only if the original graph has an
\(\ell\)-clique. Let \(n\) denote the number of vertices in this new graph.

Apply Lin's reduction~\cite[Lemma~14]{lin} to this \(\ell\)-partite graph.
In polynomial time it produces a family of \(m\le n^2\) sets on a universe
of size \(u\le \ell^3\lceil\log n\rceil\). The graph contains an
\(\ell\)-clique if and only if at most \(b=\binom{\ell}{2}\) of the sets
cover the universe. We can answer yes immediately if the universe is empty,
and no if some element belongs to no set.

We adapt the binary address construction of Froese~\cite[Theorem~4.1]{froese}.
Number the elements from \(1\) to \(u\), and let
\(h=\lceil\log(u+1)\rceil\). We construct a matrix with \(h\) address
columns followed by one column for each set. Include one all-zero row.
For each element \(i\), let \(a_i\) be its \(h\)-bit binary representation
and let \(v_i\) be its \(m\)-bit incidence vector in the set family.
Add the two rows \((a_i,0^m)\) and \((a_i,v_i)\), and set \(k=h+b\).
Different elements have different nonzero addresses, and each \(v_i\) is
nonzero. Hence the full rows are pairwise distinct.

We show that a feasible selection of columns exists if and only if there
is a cover of size at most \(b\). For each \(0\le j<h\), the universe
contains the element \(2^j\), since \(2^{h-1}\le u\). Its row with zero
set entries differs from the all-zero row only in address column \(j\),
numbering bits from zero. Thus every feasible selection contains all \(h\)
address columns. To separate the two rows of each element, the remaining at
most \(b\) columns must include a set containing that element. Conversely,
select all address columns and the columns of a cover of size at most \(b\).
The address columns distinguish rows of different addresses, and the cover
distinguishes the two rows at each nonzero address.

It remains to bound the running time. The matrix and its budget have
polynomial length in \(n\). Since \(n\ge 2\) and
\(u+1\le 3\ell^3\log n\), the budget satisfies
\[
 k\le g(\ell)+\log\log n,\qquad
 g(\ell)=\binom{\ell}{2}+\lceil\log(3\ell^3)\rceil+1.
\]
Fix a rational \(\rho\) with \(\gamma<\rho<1\). Writing \(x=\log n\),
the exponent hidden in the outer power of two is at most
\(A(\ell)x^\rho\) for a computable integer function \(A\ge 1\):
we can use an integer constant bounding the running-time coefficient
times \(2^{g(\ell)}\). Choose an integer \(s\ge 1/(1-\rho)\).
If \(x\ge A(\ell)^s\), then \(A(\ell)x^\rho\le x\).
Otherwise \(A(\ell)x^\rho\le A(\ell)^{s+1}\).
Thus the running time is at most
\(2^{A(\ell)^{s+1}}n^{\mathcal O(1)}\). Since \(n\) is \(\ell\) times
the original number of vertices, this decides \problem{Clique} in
fixed-parameter time, a contradiction.
\end{proof}

\begin{thebibliography}{3}
\bibitem{df} R.~G. Downey and M.~R. Fellows.
Fixed-parameter tractability and completeness II: On completeness for W[1].
\emph{Theoretical Computer Science} 141(1--2), 109--131, 1995.
\href{https://doi.org/10.1016/0304-3975(94)00097-3}{doi:10.1016/0304-3975(94)00097-3}.
\bibitem{froese} V.~Froese. \emph{Combinatorial Feature Selection:
Parameterized Algorithms and Complexity}. Master's thesis, Technische
Universit\"at Berlin, 2012, Section~4.1.1, pp.~33--34.
\href{https://fpt.akt.tu-berlin.de/publications/theses/MA-vincent-froese.pdf}{Full text}.
\bibitem{lin} B.~Lin. A simple gap-producing reduction for the parameterized
Set Cover problem. In \emph{ICALP 2019}, LIPIcs 132, 81:1--81:15.
\href{https://doi.org/10.4230/LIPIcs.ICALP.2019.81}{doi:10.4230/LIPIcs.ICALP.2019.81}.
\end{thebibliography}
\end{document}
